\documentclass[11pt]{article}
\usepackage{amsmath,amssymb,amsthm}
\usepackage[margin=1.1in]{geometry}
\usepackage[hidelinks]{hyperref}
\hypersetup{pdftitle={Parity families and signed spectra: kernel averaging, near-Ramanujan bounds, and exact circulant models}, pdfauthor={Vaibhav Suvagiya}, pdfsubject={Signed graph spectra, parity families, and near-Ramanujan bounds}, pdfkeywords={Bilu-Linial conjecture, signed graphs, parity families, Ihara zeta, non-backtracking walks, circulants}}
\setlength{\emergencystretch}{3em}

\newtheorem{lemma}{Lemma}
\newtheorem{proposition}[lemma]{Proposition}
\newtheorem{theorem}[lemma]{Theorem}
\newtheorem{corollary}[lemma]{Corollary}
\newtheorem{conjecture}[lemma]{Conjecture}
\theoremstyle{remark}
\newtheorem{remark}[lemma]{Remark}
\newtheorem{question}[lemma]{Question}

\newcommand{\F}{\mathbb{F}}
\newcommand{\R}{\mathbb{R}}
\newcommand{\Ze}{Z_{\mathrm{ev}}}
\newcommand{\rhomax}{\rho}
\DeclareMathOperator{\tr}{tr}

\title{Parity families and signed spectra: kernel averaging,\\ near-Ramanujan bounds, and exact circulant models}
\author{Vaibhav Suvagiya\\
\small Sardar Vallabhbhai National Institute of Technology, Surat}
\date{September 2026}

\begin{document}
\maketitle

\begin{abstract}
We develop an affine $\mathbb F_2$ framework for structured signings of regular graphs.
A family-averaging identity converts even spectral moments into parity-weighted
closed-walk counts supported on the span of prescribed short even cycles, while a
kernel-averaged Ihara identity gives the corresponding decomposition at the
non-backtracking level. We give a finite-scale bounded-rank counting estimate and a
conditioning corollary showing that, on bicycle-free graph sequences, any parity family
of uniformly bounded codimension contains near-Ramanujan signings whenever the
corresponding random-signing theorem applies. The latter is a transfer statement rather
than a new concentration theorem. Finally, on $C_n(1,2)$ for even $n\ge10$, the
quadrilateral-unbalanced family has exactly four switching classes and its twisted
classes attain
$\rho_-(n)=2\sqrt{\cos^2(\pi/n)+\cos^2(2\pi/n)}$; a period-$8$ signing has
spectral radius $r_*=2.793604493334841\ldots$ for every positive multiple of $8$.
Thus for $n=8m\ge32$ the constrained minimum is strictly larger than a value
attained by an unrestricted signing, while equality of $r_*$ with the unrestricted
minimum remains conjectural.
\end{abstract}

\section{Introduction}

Bilu and Linial \cite{bilulinial06} asked whether every $d$-regular
graph admits a signing $\sigma\colon E\to\{\pm1\}$ of its adjacency
matrix with $\rho(A_\sigma)\le2\sqrt{d-1}$, the spectral radius of the
infinite $d$-regular tree. Since the new eigenvalues of a $2$-lift are
the eigenvalues of a signed adjacency matrix, such a signing would
produce Ramanujan $2$-lifts. Marcus, Spielman and Srivastava \cite{mss15}
proved the one-sided bound $\lambda_{\max}(A_\sigma)\le2\sqrt{d-1}$ via
interlacing families, resolving the bipartite case. Mohanty, O'Donnell
and Paredes \cite{mop20} showed that a uniformly random signing is
near-Ramanujan when every ball of radius $r\gg(\log\log n)^2$ contains
at most one cycle. Recent general two-sided bounds are reported in the preprints of Huang
\cite{huang26}, with $2\sqrt{2(d-1)}$, and Lin and Zhou \cite{linzhou26}, with
$\frac{3+\sqrt5}{2}\sqrt{d-1}$. A September 2026 preprint by Xu \cite{xu26counter} constructs a finite connected cubic
graph for which every signing has an eigenvalue outside
$[-2\sqrt2,2\sqrt2]$; the preprint therefore reports a counterexample
to the Bilu-Linial conjecture in its full generality. The results below therefore concern structured parity
families, finite-scale counting, and exact models rather than the unrestricted statement.

This paper develops an affine $\mathbb F_2$ method for structured signings. The spectrum of
$A_\sigma$ depends only on the switching class, and the master averaging identity
(Proposition~\ref{prop:master}) expresses the family-averaged even trace as a signed sum
of parity classes in the constraint span $W$. The associated kernel-averaged Ihara identity
(Proposition~\ref{prop:Lident}) gives the same decomposition at the non-backtracking
level. These identities isolate the cancellation problem: walks whose parity lies outside
$W$ average away, while single, pair, and higher-order confined classes enter with alternating
parity weights.

The quantitative results have a deliberately narrower scope. Proposition~\ref{prop:upper}
is a finite-scale bounded-rank counting estimate based on maximal fresh runs and explicit
non-backtracking closure of stale segments. A separate conditioning corollary
(Proposition~\ref{prop:bf-conditioning}) transfers the known random-signing theorem of
Mohanty--O'Donnell--Paredes to any affine parity family of uniformly bounded codimension
on a bicycle-free graph sequence. This transfer gives a near-Ramanujan signing in the
family, but does not exploit the alternating parity weights quantitatively; the genuinely
new cancellation problem therefore remains open in the cycle-dense regime. The fixed-rank
counting condition is not presented as an asymptotic hypothesis at scale $(\log n)^2$.

Finally, the paper gives an exact circulant model on $C_n(1,2)$. The quadrilateral-unbalanced
family has four switching classes and an explicitly identified twisted minimum, while an
explicit period-$8$ signing gives a strictly smaller unrestricted value for every
$n=8m\ge32$. This exact calculation provides the main fully proved finite model in which
constrained and unrestricted spectral optimization can be separated.

The examples show why an unsigned growth restriction is needed for this
counting approach: a $K_d$ subgraph traps totally even walks at rate
$d-2>\sqrt{d-1}$ (Proposition~\ref{prop:trap}). This does not claim that the
resulting hypotheses are necessary for every possible proof. Along the way we give an exact
certificate on the hypercube: every solution of the quadrilateral
system satisfies $A_\sigma^2=nI$, recovering the signed-hypercube construction of Huang
\cite{huang19} as a special case of the family
(Proposition~\ref{prop:cube}), and we record why the natural
two-sided extension of the interlacing method fails at the first
gate: $\mathbb E_\sigma\det(xI-A_\sigma^2)$ is not real-rooted, with
the exact identity $q_{C_4}=(x^2-4x+2)^2+4$
(Remark~\ref{rem:qG}). The finite algebraic identities and the reported computational
instances were independently checked by exact or symbolic computation
where applicable; these checks do not replace the general proofs. The
experiments, which guided the constructions and falsified two
intermediate conjectures,
are summarized in Section~\ref{sec:experiments}. Exact spectral
results for signed circulants arising from the same family are developed in
Section~\ref{sec:circulants}, including an infinite family of counterexamples to the
global-optimality statement suggested by the small-size computations.



\section{Related work and positioning}\label{sec:related}

The parity component of the present framework is closely related to work that
predates this manuscript. Chen, van Dam, and Bu showed that the number of
parity-closed walks of a fixed length is the corresponding spectral moment
averaged over all signings of a fixed underlying graph~\cite{chenvandambu24}.
Our use of this observation is different in scope: we average over an affine
subfamily defined by prescribed signs on short even cycles. Fourier cancellation
then leaves a signed sum over the parity classes lying in the constraint span
$W$, rather than only the totally even class that survives under the full uniform
signing average.

There is also an existing signed-graph zeta-function literature. Li, Hou, and
Wang define a zeta function for a signed graph and obtain a determinant expression
through the signed oriented line graph~\cite{lihouwang24}. Our Proposition~\ref{prop:Lident}
uses the usual non-backtracking/Ihara structure for each signing and studies its
average over an affine switching family; the resulting kernel decomposition is
the additional feature used here. The homological character viewpoint of
Luo and Roy is likewise closely related: their 2026 preprint studies twisted
adjacency operators over graph homology characters, together with trace and
weighted cycle-counting formulas~\cite{luoroy24}. We restrict to the
$\mathbb F_2$ parity quotient relevant to signings and use character averaging
to obtain existence bounds for constrained signings, rather than to study the
full character torus.

The circulant calculations in Section~\ref{sec:circulants} first appeared in the
author's companion preprint~\cite{companion}. The present paper does not claim those
finite calculations as new relative to that preprint. Instead, it integrates the exact
circulant model with the affine parity-family framework, adds the period-$8$ construction,
and makes the constrained-versus-unrestricted distinction explicit. The earlier parity-family
preprint~\cite{prior26} contains the master averaging identity and the kernel-averaged
$L$-function; the present manuscript consolidates those results, adds the exact circulant
component, and revises the quantitative discussion and verification protocol.

For the bicycle-free existence statement, Mohanty, O'Donnell, and Paredes proved that a
uniformly random signing is near-Ramanujan when radius-$r$ neighborhoods contain at most
one cycle and $r\gg(\log\log n)^2$~\cite{mop20}. Proposition~\ref{prop:bf-conditioning}
is an explicit conditioning consequence of that theorem for affine parity families of
bounded codimension. Consequently, it should not be read as a new random-signing
concentration theorem; the new quantitative question for this manuscript remains whether
the parity weights in Proposition~\ref{prop:master} can be exploited in cycle-dense regimes.

For the Bilu-Linial problem itself, the current literature has changed materially
since the original conjecture. Huang reports a $2\sqrt{2(d-1)}$ two-sided signing
bound~\cite{huang26}; Lin and Zhou report the general coefficient
$(3+\sqrt5)/2$ in front of $\sqrt{d-1}$~\cite{linzhou26}; and Xu reports a
finite connected cubic counterexample to the original unrestricted conjecture
~\cite{xu26counter}. These are used only as literature context: the present paper
makes no claim about the unrestricted conjecture beyond the exact circulant model.
This manuscript supersedes the author's earlier preprint
``Parity families and a kernel-averaged $L$-function for near-Ramanujan signings,''
arXiv:2607.17343~\cite{prior26}. The present version substantially expands
that work by adding the exact circulant analysis and periodic counterexample,
refining the confined-walk proof architecture, and consolidating the verification
and exploratory computational methodology. The earlier preprint should therefore
be regarded as a prior version of the present manuscript, not as an independent
claim of the same results.

\section{Setup}

$G=(V,E)$ is a simple connected graph, $|V|=n$, $|E|=m$, adjacency matrix
$A$. A signing is $\sigma\in\{\pm1\}^E$, encoded by the bit vector
$s\in\F_2^E$ with $\sigma_e=(-1)^{s_e}$; the signed adjacency matrix
$A_\sigma$ has entries $\sigma_{uv}A_{uv}$. We write
$\rhomax(A_\sigma)=\max_i|\lambda_i(A_\sigma)|$. The cycle space
$Z(G)\le\F_2^E$ has dimension $m-n+1$; the \emph{sign} of $z\in Z(G)$
under $\sigma$ is $(-1)^{s\cdot z}$, and a cycle $C$ is \emph{unbalanced}
if $s\cdot\chi_C=1$. Switching preserves the spectrum and acts trivially
on cycle signs; conversely the switching class of $\sigma$ is exactly the
linear functional $z\mapsto s\cdot z$ on $Z(G)$, so there are
$2^{m-n+1}$ switching classes.

For a family $\mathcal C$ of even cycles, call the constraint system
$s\cdot\chi_C=1$ for all $C\in\mathcal C$ \emph{consistent} if it has
at least one solution $s\in\F_2^E$; its solution set is then an affine
parity family, and we write $W=\operatorname{span}_{\F_2}\{\chi_C:C\in\mathcal C\}$.
A graph is \emph{$R$-bicycle-free} if every ball of radius $R$ has cycle
rank at most $1$, where the cycle rank of a graph $H$ is
$|E(H)|-|V(H)|+c(H)$ and $c(H)$ is its number of connected components
(equivalently, every such connected ball is a tree or unicyclic).
Throughout, graph distance and balls are taken in the underlying
unsigned graph.

The Bilu-Linial conjecture was the unrestricted statement
\cite{bilulinial06}; a September 2026 preprint by Xu \cite{xu26counter}
reports a finite counterexample to that unrestricted statement. The bounds most relevant to the present method
include the one-sided result of Marcus, Spielman and Srivastava
\cite{mss15}, the bicycle-free random-signing result of Mohanty, O'Donnell
and Paredes \cite{mop20}, and the later general two-sided bounds of Huang, and of Lin and Zhou
\cite{huang26,linzhou26}.

\section{Even-trace parity}

For a closed walk $W$ of length $\ell$ let $m_e(W)$ be the number of
traversals of $e$, and let $z(W)\in\F_2^E$ be the parity vector
$z(W)_e=m_e(W)\bmod 2$. Since $W$ is closed, every vertex has even total
incidence with the edge multiset of $W$, so the support of $z(W)$ has all
even degrees: $z(W)\in Z(G)$. Moreover
$|z(W)|\equiv\sum_e m_e(W)=\ell\pmod 2$, and the sign of $W$ under
$\sigma$ is $\prod_e\sigma_e^{m_e(W)}=(-1)^{s\cdot z(W)}$.

Let $\Ze=\{z\in Z(G):|z|\text{ even}\}$, the kernel of the linear map
$z\mapsto|z|\bmod2$ on $Z(G)$. If $G$ is bipartite then $\Ze=Z(G)$;
otherwise $\Ze$ has index $2$, hence dimension $m-n$.

\begin{lemma}[Even-trace parity]\label{lem:parity}
For every signing $\sigma$ and every even $\ell$,
\[
\tr\!\big(A_\sigma^{\ell}\big)
=\sum_{z\in \Ze} N_\ell(z)\,(-1)^{s\cdot z},
\]
where $N_\ell(z)\ge0$ is the number of closed $\ell$-walks in $G$ with
parity vector $z$, independent of $\sigma$.
\end{lemma}

\begin{proof}
Group the closed $\ell$-walks by $z(W)$ and use
$\mathrm{sign}(W)=(-1)^{s\cdot z(W)}$. For even $\ell$ only classes with
$|z|$ even occur, by the congruence above.
\end{proof}

\begin{corollary}\label{cor:evenchar}
$\rhomax(A_\sigma)$ is a function of the restriction of the character
$z\mapsto(-1)^{s\cdot z}$ to $\Ze$. In particular
$\rhomax(A_{-\sigma})=\rhomax(A_\sigma)$, and for non-bipartite $G$ the
map from switching classes to $\rhomax$-relevant data is $2$-to-$1$: at
most $2^{m-n}$ values of $\rhomax$ occur among the $2^{m-n+1}$ classes.
\end{corollary}

\begin{proof}
$\rhomax(A_\sigma)=\lim_{k\to\infty}\tr(A_\sigma^{2k})^{1/2k}$ since
$\rhomax^{2k}\le\tr(A_\sigma^{2k})\le n\rhomax^{2k}$, and by
Lemma~\ref{lem:parity} every even trace depends on $s$ only through
$s|_{\Ze}$. Negating $\sigma$ adds $|z|\bmod 2$ to $s\cdot z$, which
vanishes on $\Ze$.
\end{proof}

\begin{corollary}[Interference accounting]\label{cor:interference}
Let $C$ be a cycle of $G$.
If $|C|$ is even, the parity class $z=\chi_C$ contributes to even traces
with the sign $(-1)^{s\cdot\chi_C}$ of $C$ itself. If $|C|$ is odd, then
$\chi_C\notin\Ze$ and $C$ cannot occur alone in an even trace; the basic
combined classes are $z=\chi_{C_1}+\chi_{C_2}$ for pairs of odd cycles,
contributing with sign equal to the \emph{product} of the two cycle signs.
\end{corollary}

Corollary~\ref{cor:interference} identifies the sign structure but not
the magnitudes $N_\ell(z)$, so it does not by itself prove monotonicity
statements. It does, however, explain two robust empirical phenomena
(\S\ref{sec:experiments}): making all short \emph{even} cycles unbalanced
places a deterministic $(-1)$ on the terms corresponding to individual
short even-cycle parity classes, while making all short \emph{odd} cycles
unbalanced makes every odd-pair term
$(+1)(+1)$ or $(-1)(-1)$, i.e.\ sign-definite \emph{positive},
eliminating cancellation; and indeed exhaustive optima place roughly half
of all triangles unbalanced, while forcing all triangles unbalanced
raises $\rhomax$ above the random-signing mean.

\section{An F2 obstruction and the parameter beta L}

Fix $L$ and consider the linear system over $\F_2$
\begin{equation}\label{eq:sys}
s\cdot\chi_C=1\qquad\text{for every even cycle $C$ of length}\le L .
\end{equation}

\begin{lemma}[$K_4$ obstruction]\label{lem:k4}
If $G$ contains $K_4$ as a subgraph, then the system \eqref{eq:sys} with
$L\ge4$ is inconsistent: no signing of $G$ makes all quadrilaterals
unbalanced.
\end{lemma}

\begin{proof}
Let $\{a,b,c,d\}$ span a $K_4$ and let $C_1,C_2,C_3$ be its three
$4$-cycles. Every edge of the $K_4$ lies in exactly two of them, so
$\chi_{C_1}+\chi_{C_2}+\chi_{C_3}=0$ in $\F_2^E$, whence
$s\cdot\chi_{C_1}+s\cdot\chi_{C_2}+s\cdot\chi_{C_3}=0$ for every $s$,
while \eqref{eq:sys} demands the sum be $1+1+1=1$.
\end{proof}

More generally, any odd-cardinality dependent family of even short
cycles blocks \eqref{eq:sys}. Let $\mathcal C_L(G)$ denote the set of even
cycles of length at most $L$. If $\mathcal C_L(G)\ne\varnothing$, define
\[
\beta_L(G)\;=\frac{1}{|\mathcal C_L(G)|}\;\max_{s\in\F_2^E}
|\{C\in\mathcal C_L(G):s\cdot\chi_C=1\}|,
\]
the maximum simultaneously-unbalanceable fraction; if
$\mathcal C_L(G)=\varnothing$, set $\beta_L(G)=0$. Every explicit
signing certifies a lower bound on $\beta_L$, and $\beta_L\ge\tfrac12$
whenever $\mathcal C_L(G)$ is nonempty, by averaging over uniform
signings. For a chosen consistent subsystem, row reduction produces its affine
solution family; its dimension is determined by that subsystem.
Computing $\beta_L$ exactly is a MAX-XOR-SAT-type optimization problem;
we do not use any complexity-theoretic claim about this restricted
cycle-incidence family below. Our data
(\S\ref{sec:experiments}) suggest $\beta_L$ as a useful candidate parameter for the spectral
defect $\min_\sigma\rhomax(A_\sigma)-2\sqrt{d-1}$ on cycle-dense instances.

\section{The hypercube: an exact certificate}

\begin{proposition}\label{prop:cube}
On $Q_n$ the system \eqref{eq:sys} with $L=4$ is consistent, and
\emph{every} solution $\sigma$ satisfies $A_\sigma^2=nI$; consequently
$\rhomax(A_\sigma)=\sqrt n$ for the entire affine solution family.
\end{proposition}

\begin{proof}
Consistency is witnessed by Huang's signing \cite{huang19}. Let $\sigma$
be any solution. $Q_n$ is triangle-free, so
$(A_\sigma^2)_{uv}=0$ for adjacent $u,v$. If $u\ne v$ are at distance
$2$ they have exactly two common neighbours $w,w'$, and
$u\,w\,v\,w'$ is a $4$-cycle; unbalancedness gives
$\sigma_{uw}\sigma_{wv}=-\sigma_{uw'}\sigma_{w'v}$, so the two terms of
$(A_\sigma^2)_{uv}$ cancel. Pairs at distance $\ge3$ contribute nothing,
and the diagonal equals the degree $n$.
\end{proof}

This proves the zero variance of $\rhomax$ observed across sampled solutions on $Q_5,Q_6$ (all at $\sqrt5,\sqrt6$ exactly, against a Kesten
floor of $4$ and $2\sqrt5$): on $Q_n$ the local $\F_2$ condition is an
exact spectral certificate, far below $2\sqrt{d-1}$. It also isolates
what is special: the certificate needs \emph{every} $2$-path to be
completed by exactly one $4$-cycle. Graphs with this property are
covered by the two-eigenvalue framework of signed covers; the general
question is how much of the cancellation survives when only a
$\beta_L$-fraction of quadrilaterals can be unbalanced.


\section{\texorpdfstring{Failure of two-sided interlacing for $A_\sigma^2$}{Failure of two-sided interlacing for A sigma squared}}

By Godsil-Gutman \cite{godsilgutman81},
$\mathbb{E}_\sigma\det(xI-A_\sigma)=\mu_G(x)$, the matching polynomial,
which is real-rooted with roots in $[-2\sqrt{d-1},2\sqrt{d-1}]$
\cite{heilmannlieb72}; this underlies \cite{mss15}. A natural two-sided
analogue would take
$q_G(x):=\mathbb{E}_\sigma\det\!\big(xI-A_\sigma^2\big)$ and hope for
real-rootedness plus an interlacing family, which would give some
$\sigma$ with $\rhomax(A_\sigma)^2\le\operatorname{maxroot}(q_G)$.

\begin{remark}\label{rem:qG}
$q_G$ is not real-rooted in general. For $G=C_4$ the two switching
classes have $A_\sigma^2$-spectra $\{4,4,0,0\}$ and $\{2,2,2,2\}$, so
\[
q_{C_4}(x)=\tfrac12\big[x^2(x-4)^2+(x-2)^4\big]
=x^4-8x^3+20x^2-16x+8=(x^2-4x+2)^2+4,
\]
which has no real roots at all. On the finite atlas and standard examples
recorded in the accompanying repository, the same exact averaging procedure
finds non-real-rooted $q_G$ for $C_4$, $K_{3,3}$, $Q_3$, and Petersen,
while $q_{K_4}=x^4-12x^3+42x^2-52x+21$ is real-rooted. This is a finite
computational observation, not a classification; exact integer characteristic
polynomials are used before the real-root count. A broader characterization of
graphs with real-rooted $q_G$ remains open.
\end{remark}

The even-subgraph corrections that cancel in
$\mathbb{E}\det(xI-A_\sigma)$ survive squaring; this is the concrete
obstruction to extending the interlacing method two-sidedly over the
plain signing family. Any successful two-sided interlacing extension
must overcome this obstruction.

\section{Computational evidence}\label{sec:experiments}

The computations are supplementary checks and exploratory evidence; the
analytic statements below are not justified by the numerical experiments.
The verification suite and archived data are available in the accompanying
repository, at \url{https://github.com/Vaibhavs25/bilu-linial-parity}. The
verification suite uses Python with NumPy, NetworkX, and SymPy. Exact algebraic
checks use integer or symbolic arithmetic; floating-point eigensolvers are used
only for numerical spectral checks. Reproduction begins with
\texttt{python3~verify\_all.py}.

\emph{Exact finite checks.} The defect atlas exhausts connected regular
graphs on at most $7$ vertices and several standard examples, including
Petersen, $Q_3$, $Q_4$, Heawood, M\"obius-Kantor, and the circulants
$C_n(1,2)$. Switching classes are counted once, using co-tree sign
coordinates. The reported percentages refer to the stated finite graph
collection only and are not used as evidence for the asymptotic theorems.
For $Q_4$ the exhaustive optimum is consistent with the signed-hypercube
phenomenon; for $Q_5,Q_6$ the even-cycle system is satisfiable and every
solution has $\rhomax=\sqrt n$ by Proposition~\ref{prop:cube}. On even
$C_n(1,2)$ the quadrilateral system is satisfiable and the constrained
family has $\rhomax\le2\sqrt2$, with the exact values given in
Section~\ref{sec:circulants}. These calculations are exact or
switching-class exhaustive where stated.

\emph{Variance-onset exploration.} On random $2$-lift towers and
planted-quadrilateral $4$-regular graphs we separate the deterministic
row-reduction solution from the randomly sampled affine tail before
computing dispersion. The earlier zero-tail outlier is retained only as
provenance and is excluded from tail-dispersion statistics by definition.
The resulting tower ratios from $0.79$ to $1.21$ and the reported finite-size
Kesten benchmarks are empirical observations. In planted instances the
full short-even system can be inconsistent; the code therefore studies a
retained, greedy-consistent parity subsystem. No causal interpretation is
assigned to the resulting differences from uniform random signings.

\emph{Exploratory campaign.} The current repository version of
\texttt{campaign.py} uses an equal spectral-evaluation budget of $2000$
spectral-radius evaluations per strategy and records the rank of the
retained parity subsystem. Its random-affine ablation now constructs a
full-row-rank system with exactly that same rank on each instance. The
planted population and the MCMC-conditioned graph control are not claimed
to have identical $C_4$ distributions, and the latter is explicitly an
MCMC control rather than a proven uniform conditional sampler. The earlier
campaign data used a different ablation protocol and are retained as
provenance; they are not used to support a claim that an observed effect is
specifically due to parity. All statistical comparisons in this section
are therefore descriptive/exploratory. An equal spectral-evaluation budget
also does not mean equal wall-clock computational cost, which is recorded
separately. The repository archives the verification scripts and finite data
used for these checks.

\section{The kernel-averaged trace}\label{sec:kernel}

Fix a consistent system \eqref{eq:sys}, let
$W=\operatorname{span}\{\chi_C\}\subseteq\Ze$ be the span of its
constraint cycles, and let $\mathcal F$ be its family of switching
classes.

\begin{proposition}[Master identity]\label{prop:master}
The map $\pi\colon W\to\F_2$,
$\pi\big(\sum_{C\in S}\chi_C\big)=|S|\bmod2$, is a well-defined linear
form, and for every even $\ell$,
\[
\mathbb E_{\sigma\in\mathcal F}\,\tr\!\big(A_\sigma^{\ell}\big)
\;=\;\sum_{z\in W}(-1)^{\pi(z)}\,N_\ell(z).
\]
\end{proposition}

\begin{proof}
Consistency means every dependency $\sum_{C\in S}\chi_C=0$ has $|S|$
even, so $\pi$ is well defined; linearity follows from
$|S|+|S'|\equiv|S\triangle S'|\pmod2$. Identify switching classes with
$\operatorname{Hom}(Z(G),\F_2)$; then
$\mathcal F=\varphi_0+\operatorname{Ann}(W)$ and $\varphi|_W=\pi$ for
every $\varphi\in\mathcal F$. For $z\in Z(G)$ the average
$\mathbb E_{\psi\in\operatorname{Ann}(W)}(-1)^{\psi(z)}$ equals $1$ if
$z\in(\operatorname{Ann}W)^{\perp}=W$ and $0$ otherwise. Averaging the
expansion of Lemma~\ref{lem:parity} term by term gives the claim; note
$W\subseteq\Ze$ since the constraint cycles are even.
\end{proof}

Walks whose wrap parity touches any cycle outside $W$ are annihilated;
a single constraint cycle enters with $-1$, a pair with $+1$, and so on.

\begin{corollary}\label{cor:certificate}
For every $k$,
\[
\min_{\sigma\in\mathcal F}\rhomax(A_\sigma)\;\le\;
R_{\mathcal F}(k):=\Big(\sum_{z\in W}(-1)^{\pi(z)}N_{2k}(z)\Big)^{1/2k}.
\]
By contrast the average over all signings retains only $z=0$:
$\mathbb E_\sigma\tr(A_\sigma^{2k})=N_{2k}(0)\ge n\,t_{2k}(d)$, where
$t_{2k}(d)$ counts closed walks at the root of the infinite $d$-regular
tree; indeed every such tree walk crosses each tree edge an even number
of times and projects, injectively by unique lifting, to a closed walk
in $G$ with all edge multiplicities even. Hence the uniform average can
never certify a spectral radius below the Kesten profile, while
$R_{\mathcal F}(k)$ can, if and only if the signed sum over
$W\setminus\{0\}$ is negative enough to consume both the even-wrap
excess $N_{2k}(0)-n\,t_{2k}(d)$ and part of the tree term.
\end{corollary}

\begin{corollary}[Confined mass]\label{cor:balanced}
The homogeneous family $\operatorname{Ann}(W)$, all constraint cycles
\emph{balanced}, measures the total parity-confined walk mass:
\[
\mathbb E_{\psi\in\operatorname{Ann}(W)}\,\tr\!\big(A_\psi^{2k}\big)
=\sum_{z\in W}N_{2k}(z).
\]
This is the same annihilator computation with $\varphi_0=0$, so every
$z\in W$ enters with sign $+1$; it makes the confined mass, and hence
Question~\ref{q:confined} below, directly measurable by sampling
eigenvalues.
\end{corollary}

The finite computations in $\S9$ are used only descriptively. In particular, the
quantity $R_{\mathcal F}(k)$ in Corollary~\ref{cor:certificate} is an unnormalized
trace-root upper bound on $\min_\sigma\rhomax(A_\sigma)$; a normalized moment root would
differ by the factor $n^{1/(2k)}$. The numerical values reported in this section are estimated
from sampled signings and their spectra and are therefore Monte-Carlo estimates, not rigorous
spectral-radius certificates or asymptotic evidence. On the planted-quadrilateral $4$-regular
samples at fixed density $\tfrac12$ quadrilateral per vertex,
$n\in\{128,256,512\}$, $k=30$, the parity-corrected trace ratio remains below one
(approximately $0.08$--$0.11$), while the balanced-family excess stays roughly $n$-stable
(about $0.12$--$0.15$ in the corresponding moment roots). These are descriptive finite
computations obtained by sampling.

\section{\texorpdfstring{A kernel-averaged $L$-function}{A kernel-averaged L-function}}\label{sec:zeta}

By Bass's theorem \cite{bass92}, for any signing
$\det(I-uB_\sigma)=\prod_{p}\big(1-\sigma(p)u^{|p|}\big)$, the product
running over primes (primitive cyclic classes of closed non-backtracking
walks), where $\sigma(p)=(-1)^{\varphi(z(p))}$ and $z(p)\in Z(G)$ is the
edge-multiplicity parity vector of $p$. Averaging over the family
diagonalizes this product.

\begin{proposition}[Kernel-averaged $L$-identity]\label{prop:Lident}
With $W$, $\pi$, $\mathcal F$ as above, $\varepsilon_p=(-1)^{\pi(z(p))}$,
and $B$ the \emph{unsigned} non-backtracking operator, define
$\log\det(I-uB_\sigma)$ to be the analytic branch that vanishes at $u=0$.
Then, for $|u|<\tfrac1{d-1}$,
\[
\mathbb E_{\sigma\in\mathcal F}\,\log\det(I-uB_\sigma)
=\tfrac12\log\det(I-u^2B)
-\tfrac12\!\!\sum_{p\,:\,z(p)\in W}\!\!
\log\frac{1+\varepsilon_pu^{|p|}}{1-\varepsilon_pu^{|p|}}\,.
\]
\end{proposition}

\begin{proof}
For $|u|<1/(d-1)$ the Euler product and its logarithm converge
absolutely. With the branch fixed by $\log\det(I-uB_\sigma)=0$ at $u=0$,
the character $(-1)^{\varphi(z)}$ is fixed at $(-1)^{\pi(z)}$ on $\mathcal F$ when
$z\in W$, whereas for $z\notin W$ its average over the affine family is zero. Hence
$\mathbb E\log(1-\sigma(p)u^{|p|})$ equals
$\log(1-\varepsilon_pu^{|p|})$ for parity-confined primes and
$\tfrac12\log(1-u^{2|p|})$ for the rest. This uses only linearity of expectation;
no independence between the signs of different primes is assumed. Collecting the
factor over all primes gives
\emph{all} primes gives
$\tfrac12\log\prod_p(1-(u^2)^{|p|})=\tfrac12\log\det(I-u^2B)$, and the
per-prime correction for confined primes is
$\log(1-\varepsilon u^{\ell})-\tfrac12\log(1-u^{2\ell})
=\tfrac12\log\frac{1-\varepsilon u^{\ell}}{1+\varepsilon u^{\ell}}$.
Absolute convergence for $|u|(d-1)<1$ follows from
$\#\{p:|p|=\ell\}\le\tfrac{2m}{\ell}(d-1)^{\ell}$.
\end{proof}

Both sides transfer to the adjacency side member-wise via
$\det(I-uB_\sigma)=(1-u^2)^{m-n}\det(I-uA_\sigma+u^2(d-1)I)$. The
identity was verified exactly on the prism ($\mathcal F$ of size $2$,
$W$ spanned by the three quadrilaterals): the coefficient form
$\mathbb E_{\mathcal F}\tr B_\sigma^k$ matches the prime decomposition
for all $k\le14$ to machine precision, and the displayed identity holds
to $8$ digits at $u=0.2$ under truncation at $|p|\le14$.

The structural content: the untwisted term $\tfrac12\log\det(I-u^2B)$
extends analytically to the Ramanujan disk $|u|<(d-1)^{-1/2}$; primes
whose parity escapes $W$ are averaged into this untwisted contribution,
while the remaining correction is carried by parity-confined primes. At the level of the averaged $L$-function, the sign problem has
become a counting problem. We emphasize the scope of the identity: the
domain of analyticity of $\mathbb E\log\det$ is governed by the worst
member of $\mathcal F$, so it does not by itself bound
$\min_{\sigma\in\mathcal F}\rho$; the quantitative route of this
paper is Corollary~\ref{cor:certificate}, and the identity's role is
to explain the mechanism and to isolate Question~\ref{q:confined}.

\begin{question}[Confined-prime counting]\label{q:confined}
For which graph sequences does
$\#\{p:|p|=\ell,\ z(p)\in W\}\le C\,(\sqrt{d-1}+\delta)^{\ell}$ hold?
Parity confinement means the odd-multiplicity edge set of $p$ is a sum
of short even cycles. The extreme case $z=0$ (totally even primes) counts
non-backtracking closed walks with every edge multiplicity even. For a fixed
finite connected $d$-regular graph this slice has exponential rate $d-1$ in
its walk length, whereas the $\sqrt{d-1}$ scale is the lower benchmark supplied
by the doubling construction and is relevant when the underlying graphs vary.
The question is whether nonzero confined classes admit a uniform counting bound
strong enough to make the parity weights in Proposition~\ref{prop:master}
quantitatively useful. Caution from the prism: on graphs so small that $W$
nearly fills the cycle space, confinement is no restriction (measured rate
$\approx1.84$ against $\sqrt2\approx1.41$ at $d=3$).
\end{question}

\section{\texorpdfstring{The $z=0$ slice: totally even walks}{The z=0 slice: totally even walks}}\label{sec:even}

Sub-question of Question~\ref{q:confined} at $z=0$: count closed
non-backtracking walks all of whose edge multiplicities are even
(``totally even'' walks). The doubling argument supplies a universal lower benchmark, but for a fixed finite
regular graph the totally even slice actually has the full non-backtracking growth rate.
The difficulty is to obtain a uniform bound for graph sequences without dense-core hypotheses.

\begin{proposition}[Slice identity and doubling]\label{prop:slice}
Let $\mathrm{Ev}_\ell$ denote the number of totally even cyclically
non-backtracking closed walks of length $\ell$ in $G$. Then
$\mathrm{Ev}_\ell=\mathbb E_{\sigma}\,\tr(B_\sigma^{\ell})$, the average
over all signings. For every integer $q\ge1$, the doubling map
$w\mapsto w^2$ injects cyclically non-backtracking closed walks of length
$q$ into totally even walks of length $2q$, so
\[
\mathrm{Ev}_{2q}\ge\tr(B^q).
\]
Consequently,
\[
\liminf_{q\to\infty}\mathrm{Ev}_{4q}^{1/(4q)}\ge\sqrt{d-1}.
\]
\end{proposition}

\begin{proof}
$\tr(B_\sigma^\ell)$ is the signed count of cyclically non-backtracking
closed walks, the sign being $\prod_e\sigma_e^{m_e(w)}$; averaging kills every
walk with an odd multiplicity. If $w$ is cyclically non-backtracking
then so is $w^2$ (the seam repeats a legal transition of $w$), its
multiplicities double, and $w$ is recovered as the first half. For
bipartite $G$, $\tr(B^q)=0$ for odd $q$, so the subsequence $\ell=4q$
avoids this parity obstruction.
\end{proof}

\begin{remark}[Fixed-graph growth of the totally even slice]
Assume $G$ is a fixed connected $d$-regular graph with $d\ge3$, and let
$r=|E|-|V|+1$ be its cycle rank. For every even $\ell$,
\[
0\le \mathrm{Ev}_\ell\le \operatorname{tr}(B^\ell).
\]
Partition the uniform signing space into its $2^r$ switching classes. The all-plus class
has non-backtracking spectral radius $d-1$. If $G$ is non-bipartite, the all-minus signing
is a distinct switching class and satisfies $B_{-\mathbf 1}=-B$, so its trace agrees with the
all-plus trace on even powers. If $G$ is bipartite, all-minus is switching-equivalent to
all-plus, and the unsigned non-backtracking operator has the peripheral pair
$\pm(d-1)$. By the equality case of Wielandt's theorem, every switching class other than
these extremal contributions has spectral radius strictly below $d-1$; in the extremal
class(es), every non-peripheral eigenvalue is also strictly smaller in modulus. Since $G$
is fixed and there are finitely many switching classes, there is $\rho_2<d-1$ bounding all
these subleading eigenvalue moduli. Writing $N=2|E|$, the aggregate contribution of all
subleading terms is $O(N\rho_2^\ell)$ after averaging. Consequently, for even $\ell$,
\[
\mathrm{Ev}_\ell=2^{-r}\,2(d-1)^\ell+O(N\rho_2^\ell).
\]
In the bipartite case the leading factor $2(d-1)^\ell$ comes from the peripheral pair
$\pm(d-1)$ in the single extremal switching class; in the non-bipartite case it comes from
the all-plus and all-minus switching classes, each contributing $(d-1)^\ell$ on even powers.
Since $\rho_2<d-1$,
\[
\lim_{\substack{\ell\to\infty\\ \ell\equiv0\pmod2}}
\mathrm{Ev}_\ell^{1/\ell}=d-1.
\]
Thus the $\sqrt{d-1}$ value in Proposition~\ref{prop:slice} is only a doubling lower
benchmark, not the fixed-graph exponential rate, and this fixed-graph statement does not
provide a uniform asymptotic estimate for a varying graph sequence.
\end{remark}

\begin{proposition}[Dense trapping]\label{prop:trap}
Let $d\ge4$ and suppose $G$ contains $K_d$ as a subgraph (realizable in
connected $d$-regular graphs: two copies of $K_d$ joined by a perfect
matching). Then there are $c_d>0$ and $\ell_0(d)$, independent of $G$
and $n$, with $\mathrm{Ev}_\ell(G)\ge c_d\,(d-2)^{\ell}$ for all even
$\ell\ge\ell_0$. Since $d-2>\sqrt{d-1}$ for every $d\ge4$, the confined
mass exceeds the profile $n(\sqrt{d-1}\,)^{\ell}$ throughout
$\ell\ge C_d\log n$.
\end{proposition}

\begin{proof}
Here $d$ denotes the degree of the ambient regular graph $G$; the clique $K_d$ itself is
$(d-1)$-regular, so its unsigned non-backtracking operator has Perron value $d-2$.
Totally even walks of the $K_d$ copy are totally even walks of $G$, and
by Proposition~\ref{prop:slice} their number is
$2^{-|E(K_d)|}\sum_{\sigma\in\{\pm1\}^{E(K_d)}}\tr(B_{K_d,\sigma}^\ell)$.
The unsigned $B_{K_d}$ is irreducible and aperiodic with Perron value
$d-2$; its remaining eigenvalues have modulus at most
$\max(1,\sqrt{d-2})<d-2$ (Ihara-Bass from
$\operatorname{spec}A(K_d)=\{d-1,(-1)^{d-1}\}$). By the equality case of
Wielandt's theorem \cite[Thm.~8.3.11]{hornjohnson13}, a signed version satisfies
$\rho(B_{K_d,\sigma})=\rho(B_{K_d})$ if and only if
$B_{K_d,\sigma}=\pm D B_{K_d}D^{-1}$ for a diagonal $D$ with
$|D|=I$, i.e.\ iff $\sigma$ is switching-equivalent to the all-plus or
all-minus signing. Hence $\rho_2:=\max$ over the remaining classes is
strictly below $d-2$, the two trivial classes contribute
$\tr B^\ell+(-1)^\ell\tr B^\ell=2(d-2)^\ell(1+o(1))$ at even $\ell$, and
the class average is at least
$2^{-r}\big[2(d-2)^\ell(1-o(1))-2^{r}\,d(d-1)\rho_2^{\ell}\big]$ with
$r=\binom d2-d+1$.
\end{proof}

\begin{remark}[Necessity of an unsigned core restriction]
Since $0\in W$ always, Proposition~\ref{prop:trap} shows that an unconditional upper
bound on the totally even mass at the Ramanujan rate cannot hold once a fixed dense core
is allowed to contribute for lengths of order $\log n$. In particular, the obstruction is
carried by the all-plus and all-minus switching classes of the core, whose non-backtracking
spectral radius is the unsigned value $d_H-1$. This is why a successful counting hypothesis
must constrain unsigned local growth, rather than merely exclude a subset of nontrivial
signing classes. The finite atlas also illustrates the distinction between the two operators:
for $K_4$ and Petersen, nontrivial classes can have adjacency spectral radius strictly below
the Kesten adjacency benchmark. The corresponding non-backtracking statement is a separate
spectral fact and should not be conflated with the adjacency comparison.
\end{remark}

The following finite-scale counting condition is useful for the
bounded-rank part of the method. Say $G$ is \emph{$(r_0,\delta)$-subcritical
at scale $\ell$} if every connected subgraph $H$ of $G$ with minimum degree
$\ge2$ and at most $\ell$ edges satisfies: (i) the cycle rank of $H$
is at most $r_0$; and (ii) the number $c_H(t)$ of closed non-backtracking walks of length $t$
(the initial/final step need not satisfy the cyclic condition) obeys
$c_H(t)\le 2|E(H)|\,(t+1)^{3r_0}(1+\delta)^{t}$ for all $t\le8\ell$.
The notion is monotone: subcriticality at scale $\ell$ implies it
at every smaller scale, and all subsequent hypotheses are stated
through this single definition.
Condition (ii) holds, for instance, whenever every thread of $H$ (a
maximal path through degree-$2$ vertices) has length at least
$\lambda$, with $1+\delta=(2r_0-1)^{1/\lambda}$: closed walks in $H$
project to reduced cyclic words in a free group of rank $\le r_0$ whose
generators cost at least $\lambda$ steps each.

\begin{lemma}[Non-backtracking reachability]\label{lem:reach}
Let $H$ be connected with minimum degree $\ge2$, $s$ edges, and not a
cycle. Between any two directed edges of $H$ there is a
non-backtracking path of length at most $2s-1$.
\end{lemma}

\begin{proof}
Let $\overrightarrow H$ be the oriented line graph whose vertices are the
$2s$ directed edges of $H$, with an arc $e\to f$ when $t(e)=o(f)$ and
$f\ne\bar e$. For a connected graph of minimum degree at least $2$ that
is not a cycle, $\overrightarrow H$ is strongly connected \cite[Section~1]{ks00}.
Hence between any two of its vertices there is a simple directed path of
length at most $2s-1$, which is exactly the required non-backtracking path in $H$.
\end{proof}

\begin{lemma}[Fresh-run rank increase]\label{lem:fresh}
Let $w$ be a cyclically non-backtracking closed walk, and expose a chosen
initial directed edge before declaring edges fresh. Let $S$ be the support
of $w$, and let $q$ be the number of maximal nonempty fresh runs. Then
\[
q\le \operatorname{rank}(S).
\]
\end{lemma}

\begin{proof}
Expose the initial directed edge first, so the current support $T$ is nonempty and connected.
For each maximal fresh run, let $F$ be the set of edges exposed during that run and let
$T'$ be the support immediately after the run. We claim
\[
\operatorname{rank}(T')-\operatorname{rank}(T)\ge1.\tag{*}
\]
The run starts at a vertex of $T$, and every edge in $F$ is new when it is exposed.

Suppose first that the fresh trail meets $T$ again before it terminates. Take its first
return to $T$. The subtrail from its starting vertex to that first return has no internal
vertex in $T$ and consists entirely of fresh edges. Adding this nonempty path to the
connected graph $T$ creates a cycle, so the cycle rank increases by at least one. The same
conclusion holds when the fresh run itself terminates in $T$.

Now suppose that the fresh run terminates at a vertex outside $T$. It cannot be the final
run of the closed walk, because the final vertex is the initial vertex, which belongs to $T$.
Hence the next step is stale. Its initial vertex lies outside $T$, while every edge of $T$ has
both endpoints in $T$, so this stale edge cannot belong to $T$. It must therefore be an edge
already exposed earlier in the same fresh run. Choose the first such reuse. Before that reuse
all edges exposed by the run are distinct. Because the walk is non-backtracking, the reused
edge is not the immediately preceding edge; the segment between its two occurrences therefore
contains a nontrivial closed trail using only edges exposed by this run. Its support contains
a cycle, and again the cycle rank increases by at least one, proving $(*)$.

Different maximal fresh runs use disjoint sets of fresh edges. Applying $(*)$ successively
and telescoping the rank increments gives
\[
\operatorname{rank}(S)-\operatorname{rank}(T_0)\ge q,
\]
where $T_0$ is the one-edge initial support and has rank $0$. Hence
$q\le\operatorname{rank}(S)$.
\end{proof}

\begin{lemma}[Closing a stale segment]\label{lem:stale}
Let $S$ be connected with minimum degree at least $2$, $s$ edges, and not
a cycle. Let $P=(e_1,\ldots,e_t)$ be a non-backtracking path in $S$.
Then, for some $a\in\{0,\ldots,2s-1\}$, $P$ is the initial segment of a
closed non-backtracking walk of length $t+a$. Consequently, for fixed $t$,
the number of such paths is at most
\[
2s\sum_{a=0}^{2s-1}c_S(t+a),
\]
where $c_S(r)$ denotes the number of closed non-backtracking walks of
length $r$ in $S$.
\end{lemma}

\begin{proof}
Write $e_1=(v_0,v_1)$ and $e_t=(v_{t-1},v_t)$. Since $S$ has minimum degree at least
$2$, choose a directed edge $h$ entering $v_0$ with $h\ne\bar e_1$. By
Lemma~\ref{lem:reach}, there is a directed path in the oriented line graph
\[
e_t=e^{(0)},e^{(1)},\ldots,e^{(a)}=h,
\qquad 0\le a\le2s-1.
\]
Every transition along this path is non-backtracking. Appending
$e^{(1)},\ldots,e^{(a)}$ to $P$ therefore preserves non-backtracking, and the final edge
$h\ne\bar e_1$ makes the cyclic seam back to $e_1$ legal. We obtain a closed
non-backtracking walk of length $t+a$ whose first $t$ edges are exactly $P$.

To count such paths, choose once and for all a shortest admissible line-graph completion for
each terminal directed edge. For fixed $t$, $a$, and initial directed edge $e_1$, the first
$t$ edges of the completed closed walk recover $P$ uniquely. Thus the completion map is
injective. There are at most $2s$ choices of the initial directed edge and
$0\le a\le2s-1$, giving
\[
\#\{P\}\le2s\sum_{a=0}^{2s-1}c_S(t+a).
\]
The outside factor $2s$ is deliberately crude: $c_S(t+a)$ already counts closed walks with
all possible starting directed edges. Retaining this harmless overcount gives a uniform bound
that is convenient in Proposition~\ref{prop:upper}.
\end{proof}

\begin{proposition}[Even-walk upper bound, bounded-rank regime]\label{prop:upper}
If $G$ is $d$-regular and $(r_0,\delta)$-subcritical at scale $\ell$,
then
\[
\mathrm{Ev}_\ell(G)\;\le\;
nd\,(C\ell)^{\,4(r_0+1)^2}\,
\Big((1+\delta)^{2(r_0+2)}\sqrt{d-1}\Big)^{\ell}
\]
for an absolute constant $C$. The statement is intended only at the stated
finite scale; no asymptotic conclusion in $\ell$ is drawn from it.
\end{proposition}


\begin{proof}
Fix the starting directed edge ($\le nd$ choices) and expose the walk
step by step; let $S_i$ be the subgraph of edges used through step $i$
and $S=S_\ell$ the support. Since multiplicities are even,
$s:=|E(S)|\le\ell/2$; a degree-$1$ vertex of $S$ would force a
backtrack, so $S$ has minimum degree $2$. By hypothesis (i),
$r:=\operatorname{rank}(S)\le r_0$.

By Lemma~\ref{lem:fresh}, the number $q$ of maximal fresh runs satisfies
$q\le r\le r_0$. These runs are separated by at most $q+1\le r_0+1$ nonempty stale
segments. For fixed $q$, specifying the beginning and end of each fresh run uses at most
$2q$ time indices, so all run schedules contribute at most $C_{r_0}\ell^{2r_0}$ possibilities.
Once the incoming directed edge is fixed, each fresh edge has at most $d-1$ choices, giving
$(d-1)^s$.

For a stale segment of length $t$, if $S$ is not a single cycle, Lemma~\ref{lem:stale} gives
\[
2s\sum_{a=0}^{2s-1}c_S(t+a)
\le C s^2(2\ell+1)^{3r_0+1}(1+\delta)^{t+2s-1},
\]
because $t+a\le2\ell-1$. If $S$ is a single cycle, then after the initial directed edge is
fixed the non-backtracking continuation is forced, and the contribution is at most $2s$, which
is absorbed by the same polynomial factor. There are at most $r_0+1$ stale segments, and their
total length is exactly the number of non-fresh steps,
\[
(\ell-1)-(s-1)=\ell-s.
\]
Multiplying the segment bounds therefore gives
\[
\mathrm{Ev}_\ell\le nd\,\ell^{2r_0}
\sum_{s\le\ell/2}
\big(C s^2(2\ell+1)^{3r_0+1}\big)^{r_0+1}
(d-1)^s(1+\delta)^{(\ell-s)+(2s-1)(r_0+1)}.
\]
The exponent of $1+\delta$ is
\[
E(s)=(\ell-s)+(2s-1)(r_0+1)
=\ell+s(2r_0+1)-(r_0+1).
\]
It is increasing in $s$, and for $s\le\ell/2$,
\[
E(s)\le\ell+\frac{\ell}{2}(2r_0+1)<2(r_0+2)\ell.
\]
Also $(d-1)^s\le(d-1)^{\ell/2}$. Finally, using $s\le\ell/2$ and the at most $\ell/2$
possible values of $s$, the polynomial exponent is bounded by
\[
2r_0+2(r_0+1)+(3r_0+1)(r_0+1)+1
=(3r_0+2)(r_0+2)\le4(r_0+1)^2.
\]
After enlarging the absolute constant $C$, all remaining factors are absorbed into
$(C\ell)^{4(r_0+1)^2}$. This proves the displayed bound.
\end{proof}


\begin{remark}[Scope of the bounded-rank regime]\label{rem:scope}
The subcritical definition and Proposition~\ref{prop:upper} are used only as finite-scale
bounded-rank counting statements. They are not intended as an asymptotic fixed-$r_0$
regime at the scale $\ell\asymp(\log n)^2$. Indeed, for fixed $d\ge3$ and fixed $r$,
an elementary BFS-surplus argument finds, in every sufficiently large connected
$d$-regular graph, a connected minimum-degree-$2$ subgraph of cycle rank $r$ using
$O_{d,r}(\log n)$ edges. Thus the bounded-rank condition with fixed $r_0$ eventually
fails at every sufficiently large multiple of $\log n$, and in particular cannot support
an asymptotic theorem at the $\log^2 n$ scale. We record Proposition~\ref{prop:upper}
for the finite-scale encoding mechanism only.
\end{remark}

\begin{proposition}[Conditioning on an affine parity family]\label{prop:bf-conditioning}
Fix $d\ge3$ and $\varepsilon>0$. Let $(G_n)$ be a sequence of $d$-regular graphs on
$n$ vertices such that there are radii $R_n$ with
$R_n/(\log\log n)^2\to\infty$ and every radius-$R_n$ ball has cycle rank at most one.
Let $\mathcal F_n$ be any consistent affine parity family whose constraint span has
$\dim W_n\le K$ for a fixed constant $K$. Then, for all sufficiently large $n$, the family
contains a signing $\sigma$ with
\[
\rho(A_\sigma)\le2\sqrt{d-1}+\varepsilon.
\]
If, in addition, every prescribed cycle has length $<R_n$ and
$\Lambda_{W_n}:=\max_{z\in W_n}|z|\le\Lambda$ uniformly, then the prescribed
constraint cycles are pairwise edge-disjoint and
$\dim W_n\le\Lambda/4$.
This is a conditioning corollary of the random-signing theorem, not a new
concentration theorem.
\end{proposition}

\begin{proof}
Let $\widetilde{\mathcal F}_n\subseteq\{\pm1\}^{E(G_n)}$ be the full preimage of the
switching-class family $\mathcal F_n$. The affine cycle-sign constraints have rank
$k_n=\dim W_n$, so $\widetilde{\mathcal F}_n$ is a union of switching cosets and has
relative density
\[
p_n:=\Pr_{\sigma\,\mathrm{uniform}}(\sigma\in\widetilde{\mathcal F}_n)=2^{-k_n}\ge2^{-K}.
\]
Equivalently, the family has codimension $k_n$ in the cycle-sign quotient; the same
relative density is obtained in the full edge-sign space because every switching class has
the same number of edge-sign representatives. By the theorem of Mohanty, O'Donnell, and
Paredes~\cite{mop20}, under the stated bicycle-free hypothesis the probability of the event
$\rho(A_\sigma)>2\sqrt{d-1}+\varepsilon$ tends to zero. Hence, for all sufficiently
large $n$, the bad event has probability strictly smaller than $p_n$ and cannot
contain the whole affine family. Thus some signing in $\mathcal F_n$ satisfies the
required bound.

For the final assertion, suppose two distinct prescribed cycles shared a vertex. Since each
has length $<R_n$, both would lie in the same radius-$R_n$ ball, which would then contain two
independent cycle vectors, contradicting the bicycle-free hypothesis. Thus the prescribed
cycles are pairwise edge-disjoint. Their incidence vectors are therefore linearly independent,
so their number is $k_n=\dim W_n$. Now take the sum of all $k_n$ prescribed cycle vectors:
\[
z_*:=\sum_{i=1}^{k_n}\chi_{C_i}\in W_n.
\]
Because the cycles are edge-disjoint and every cycle has length at least $4$,
\[
|z_*|=\sum_{i=1}^{k_n}|C_i|\ge4k_n.
\]
Since $|z_*|\le\Lambda_{W_n}\le\Lambda$, we obtain
$\dim W_n=k_n\le\Lambda/4$.

\end{proof}

\begin{remark}[Scope relative to the reported cubic counterexample]
The finite cubic counterexample reported by Xu~\cite{xu26counter} is outside the
bicycle-free hypothesis here. In the construction described there, each completed
triple of leaves is joined to two new vertices adjacent to all three leaves, producing
a copy of $K_{2,3}$; this subgraph has cycle rank $2$ and diameter $2$. Thus a radius-$2$
ball already contains two independent cycles. The reported counterexample therefore does
not test the conditioning corollary above, whose hypothesis requires bicycle-freeness at
radii tending to infinity.
\end{remark}

\section{Signed circulants: exact parity families and periodic counterexamples to constrained optimality}\label{sec:circulants}

The circulant $C_n(1,2)$ has vertex set $\mathbb Z_n$ and edges
$\{i,i{+}1\}$ (step $1$) and $\{i,i{+}2\}$ (step $2$); it is
$4$-regular, with Kesten bound $2\sqrt3\approx3.464$. Its
quadrilaterals, for $n\ge10$, are exactly
$Q_i=(i,\,i{+}1,\,i{+}3,\,i{+}2)$, one for each $i\in\mathbb Z_n$, and
its triangles are $T_i=(i,\,i{+}1,\,i{+}2)$. At $n=8$ the step-$2$
edges close into two further quadrilaterals, $(0,2,4,6)$ and
$(1,3,5,7)$; we assume $n\ge10$ throughout and treat that case in
Remark~\ref{rem:n8}.

\subsection{The quadrilateral parity family}

The calculations in this section first appeared in the companion preprint
\cite{companion}. They are reproduced here because they provide the exact
finite model in which the general parity-family mechanism can be tested;
the present paper adds the period-$8$ construction and its comparison with
the constrained minimum.

We study the $\F_2$ system
\begin{equation}\label{eq:circ-system}
s\cdot\chi_{Q_i}=1\qquad(i\in\mathbb Z_n),
\end{equation}
requiring every quadrilateral to be unbalanced. This is the parity
family of the present paper, specialized to
$C_n(1,2)$; there, averaging over such families is developed into a
general method for the Bilu-Linial conjecture, with short even cycles
controlling the corresponding parity classes in the even traces. The
present graphs furnish an exactly solvable case.


\begin{proposition}\label{prop:circ}
Let $n\ge10$ be even. On $C_n(1,2)$ the system \eqref{eq:circ-system}
is consistent and its solutions form exactly four switching classes. The
class containing the signing $\sigma\equiv+1$ on step-$1$ edges
$\{i,i{+}1\}$ and $\sigma_{\{i,i+2\}}=(-1)^i$ on step-$2$ edges has
spectrum
\[
\Big\{\pm2\sqrt{\cos^2\theta_k+\cos^2 2\theta_k}\ :\
\theta_k=\tfrac{2\pi k}{n},\ 0\le k<\tfrac n2\Big\},
\]
so its spectral radius equals $2\sqrt2$ exactly, against the Kesten floor
$2\sqrt3$. For odd $n$ the system is inconsistent, and at most $n-1$
quadrilaterals can be simultaneously unbalanced.
\end{proposition}

\begin{proof}
The quadrilaterals of $C_n(1,2)$ are the $Q_i$, one per
vertex, with sign $a_ia_{i+2}\,b_ib_{i+1}$, where $a_i,b_i$ denote the
signs of $\{i,i{+}1\}$ and $\{i,i{+}2\}$. The displayed signing gives
$a_ia_{i+2}b_ib_{i+1}=(-1)^i(-1)^{i+1}=-1$ for all $i$, so the system is
consistent. If $\sum_{i\in S}\chi_{Q_i}=0$, then comparing coefficients
of the step-$2$ edge $\{i,i{+}2\}\in Q_i\cap Q_{i-1}$ forces $S$ to be
shift-invariant, so $S=\emptyset$ or $S=\mathbb{Z}_n$; the full sum is
indeed zero (each edge lies in exactly two quadrilaterals). Hence the
constraint map on the $2^{\,m-n+1}=2^{\,n+1}$ switching classes has rank
$n-1$, and the solution set is a coset of dimension $2$: four classes.

For the displayed class, write $C_1,C_2$ for the unsigned step-$1$ and
step-$2$ circulants and $D=\operatorname{diag}((-1)^j)$; then
$A_\sigma=C_1+DC_2$, which is symmetric since $D$ commutes with the even
shift $C_2$. On the orthonormal Fourier basis
$f_k(j)=n^{-1/2}e^{2\pi ijk/n}$ we have $C_1f_k=2\cos\theta_k\,f_k$,
$C_2f_k=2\cos2\theta_k\,f_k$, and $Df_k=f_{k+n/2}$. Each plane
$\operatorname{span}\{f_k,f_{k+n/2}\}$ is therefore invariant, carrying
the block
\[
\begin{pmatrix} 2\cos\theta_k & 2\cos2\theta_k\\
2\cos2\theta_k & -2\cos\theta_k\end{pmatrix},
\]
where we used $\cos\theta_{k+n/2}=-\cos\theta_k$ and
$\cos2\theta_{k+n/2}=\cos2\theta_k$. A traceless symmetric $2\times2$
block $\left(\begin{smallmatrix}a&b\\ b&-a\end{smallmatrix}\right)$ has
eigenvalues $\pm\sqrt{a^2+b^2}$, giving the stated spectrum. Setting
$u=\cos^2\theta\in[0,1]$,
$\cos^2\theta+\cos^22\theta=4u^2-3u+1\le2$ with equality iff $u=1$,
attained at $k=0$; hence $\rhomax=2\sqrt2$. For odd $n$, summing all $n$
constraints gives $0$ on the left (each edge appears twice) and
$n\equiv1\pmod2$ on the right.
\end{proof}

\begin{lemma}[Holonomy determines diagonal-unitary gauge]\label{lem:holonomy-gauge}
Let $M$ and $N$ be Hermitian matrices with the same connected support graph,
all corresponding nonzero entries of modulus one, and identical holonomy on
every cycle. Then $M$ and $N$ are conjugate by a diagonal unitary.
\end{lemma}

\begin{proof}
Choose a spanning tree and set $u_r=1$ at a root. Define $u_v$ recursively so
that $M_{uv}/N_{uv}=u_u\overline{u_v}$ on every tree edge. Equality of
holonomies on the fundamental cycles gives the same identity on every
non-tree edge. Thus $U=\operatorname{diag}(u_v)$ satisfies $M=UNU^{-1}$.
\end{proof}

\begin{proposition}[Twisted classes]\label{prop:twist}
Let $\tau_i$ denote the
sign of the triangle $T_i=(i,i{+}1,i{+}2)$ and $\alpha$ the sign product
around the step-$1$ Hamilton cycle $H$. Then
$\chi_{Q_i}=\chi_{T_i}+\chi_{T_{i+1}}$, so \eqref{eq:circ-system} is equivalent
to the alternation $\tau_{i+1}=-\tau_i$; the triangles together with $H$
form a basis of the cycle space, so the four solution classes are
coordinatized by $(\tau_0,\alpha)\in\{\pm1\}^2$; and rotation by one
vertex is an automorphism carrying $(\tau_0,\alpha)$ to
$(-\tau_0,\alpha)$. Consequently $\rhomax$ depends on $\alpha$ alone: it
equals $2\sqrt2$ for $\alpha=+1$, and for $\alpha=-1$ it equals
\[
\rho_-(n)\;:=\;2\sqrt{\cos^2(\pi/n)+\cos^2(2\pi/n)}\;<\;2\sqrt2 .
\]
In particular $\min_\sigma\rhomax(A_\sigma)\le\rho_-(n)$ for every even
$n\ge10$.
\end{proposition}

Thus $\rho_-(n)$ is the minimum spectral radius within the
quadrilateral-unbalanced family. The question of whether it is the
minimum over all signings is separate.

\begin{proof}
Each step-$2$ edge $b_i$ lies in $T_i$ only, so
$\sum_{i\in S}\chi_{T_i}=0$ forces $S=\emptyset$; together with the fact
that no combination of triangles has empty step-$2$ support except the
empty one, $\{\chi_{T_0},\dots,\chi_{T_{n-1}},\chi_H\}$ is a basis of
the $(n{+}1)$-dimensional cycle space, giving the coordinates; the
rotation statement is immediate since rotation shifts the alternating
pattern by one. For $\alpha=-1$ it therefore suffices to treat the
representative with $\tau_i=(-1)^i$. Let $R$ be the cyclic shift,
$D=\operatorname{diag}((-1)^j)$, $\varphi=\pi/n$, and set
\[
A(\varphi)\;=\;e^{i\varphi}R+e^{-i\varphi}R^{*}
\;+\;e^{2i\varphi}DR^{2}+e^{-2i\varphi}R^{-2}D ,
\]
a Hermitian matrix supported on $C_n(1,2)$. Its holonomy on the directed
triangle $i\to i{+}1\to i{+}2\to i$ is
$e^{i\varphi}\cdot e^{i\varphi}\cdot(-1)^ie^{-2i\varphi}=(-1)^i$, and on
$H$ it is $e^{in\varphi}=-1$.

The triangle cycles together with the Hamilton cycle generate the
integer cycle lattice: every step-$2$ edge occurs in exactly one triangle,
so subtracting suitable triangle cycles removes all step-$2$ coefficients
and leaves an integer multiple of $H$. Thus agreement of the triangle and
$H$ holonomies implies agreement on every integer cycle, as required by Lemma~\ref{lem:holonomy-gauge}. Hence $A(\varphi)=U A_{(+1,-1)}U^{-1}$, so the two matrices are
isospectral. Note that the step-$2$ phase is forced to be $2\varphi$: this is
exactly the co-shift that keeps the triangle holonomies real, and it is the step
omitted by a naive one-parameter Bloch argument. As in Proposition~\ref{prop:circ}, the
planes $\operatorname{span}\{f_k,f_{k+n/2}\}$ are invariant, with blocks
\[
\begin{pmatrix}
2\cos(\theta_k{+}\varphi)&2\cos(2\theta_k{+}2\varphi)\\
2\cos(2\theta_k{+}2\varphi)&-2\cos(\theta_k{+}\varphi)
\end{pmatrix},
\]
hence spectrum $\pm2\sqrt{g(\theta_k+\pi/n)}$ with
$g(\theta)=\cos^2\theta+\cos^22\theta$ and shifted momenta
$\theta_k+\pi/n\in\{(2k{+}1)\pi/n\}$, which avoid $0$ and $\pi$. Writing
$g=4u^2-3u+1$ with $u=\cos^2\theta$, $g$ decreases on $[0,\theta_0]$ and increases on
$[\theta_0,\pi/2]$, where $\cos2\theta_0=-\tfrac14$, with
$g(\pi/2)=1$ and $g$ symmetric under $\theta\mapsto\pi-\theta$;
since $g(\pi/n)\ge g(\pi/8)=\tfrac{4+\sqrt2}{4}>1$ for every even
$n\ge8$, the maximum over the shifted lattice is attained at $\pi/n$,
giving $\rho_-(n)$. (We state the inequality from $n=8$ so that
Remark~\ref{rem:n8} may invoke it.)
\end{proof}

\subsection{Periodic counterexamples}

\begin{theorem}[A periodic counterexample to constrained-family global optimality]\label{thm:counterexample}
Let $n=8m$ with $m\ge4$. On $C_n(1,2)$, assign $+1$ to every
step-$1$ edge and assign to the step-$2$ edges the period-$8$ pattern
\[
(+,+,-,+,-,-,+,-),
\]
repeated $m$ times. Let $r_*$ be the largest real root of
\[
f(x)=x^4-2x^3-6x^2+12x-4.
\]
Then the resulting signed adjacency matrix satisfies
\[
\rhomax(A_{\sigma^{(8)}})=r_*
=2.793604493334841\ldots .
\]
Moreover,
\[
r_*<\rho_-(n)
=2\sqrt{\cos^2(\pi/n)+\cos^2(2\pi/n)}
\]
for every $n=8m\ge32$. Hence
\[
\boxed{\min_\sigma\rhomax(A_\sigma)<\rho_-(n)}
\]
for every $n=32,40,48,56,\ldots$.
\end{theorem}

The signing in Theorem~\ref{thm:counterexample} is not a member of the
quadrilateral-unbalanced family. Thus the theorem does not contradict
Propositions~\ref{prop:circ} and~\ref{prop:twist}; it shows instead that
the minimizer of the constrained parity family need not be the minimizer
in the full signing space.

\begin{proof}[Proof of Theorem~\ref{thm:counterexample}]
Because the signing has period $8$, write $n=8m$ and group the vertices
into $m$ cells of eight consecutive vertices. A Bloch transform
therefore decomposes the signed adjacency matrix into $8\times8$
Hermitian blocks $H(z)$ indexed by the $m$th roots of unity $z$:
\[
H(z)=
\begin{pmatrix}
0&1&1&0&0&0&z^{-1}&z^{-1}\\
1&0&1&1&0&0&0&-z^{-1}\\
1&1&0&1&-1&0&0&0\\
0&1&1&0&1&1&0&0\\
0&0&-1&1&0&1&-1&0\\
0&0&0&1&1&0&1&-1\\
z&0&0&0&-1&1&0&1\\
z&-z&0&0&0&-1&1&0
\end{pmatrix}.
\]
Put $s=z+z^{-1}$. Since $|z|=1$, we have $s\in[-2,2]$, and a direct
calculation gives
\begin{equation}\label{eq:bloch-det}
\begin{aligned}
P(x,s)&:=\det(xI-H(z))\\
&=x^8-16x^6+80x^4-128x^2+38\\
&\qquad+s(-2x^4+16x^2-13)+s^2.
\end{aligned}
\end{equation}
At $s=2$, the determinant in \eqref{eq:bloch-det} factors as
\begin{equation}\label{eq:bloch-factor}
P(x,2)=
\bigl(x^4-2x^3-6x^2+12x-4\bigr)
\bigl(x^4+2x^3-6x^2-12x-4\bigr).
\end{equation}
The determinant formula in \eqref{eq:bloch-det} and the factorization in \eqref{eq:bloch-factor} are independently
checked by the repository's standalone exact SymPy script \texttt{verify\_theorem26.py}; running
\texttt{python3 verify\_theorem26.py} recomputes \(\det(xI-H(z))\) symbolically and verifies both identities exactly.
The verification script is pinned in the reproducibility snapshot via
\\href{https://github.com/Vaibhavs25/bilu-linial-parity/blob/b37002b020edd5d26065a0ee6ee7cff85cf7814c/verify_theorem26.py}{\\texttt{verify\\_theorem26.py}}.
Let $f(x)=x^4-2x^3-6x^2+12x-4$. Since
$f(\sqrt6)=-4$ and
\[
f'(x)=4x^3-6x^2-12x+12,
\qquad
f''(x)=12(x^2-x-1)>0\quad(x\ge\sqrt6),
\]
we have $f'(\sqrt6)=12\sqrt6-24>0$, so $f$ is strictly increasing on
$[\sqrt6,\infty)$ and has a unique root there. Direct evaluation gives
$f(2.79360)<0<f(2.79361)$; denote this root by $r_*$. For $x\ge r_*$,
\[
\frac{\partial P}{\partial s}=2s-2x^4+16x^2-13
\le -2x^4+16x^2-9<0,
\]
so $P(x,s)$ is decreasing in $s$ and hence
$P(x,s)\ge P(x,2)$ for all $s\in[-2,2]$. For $x>r_*$ the first
factor in~\eqref{eq:bloch-factor} is positive, while the difference between the second and
first factors is
\[
4x(x^2-6)>0.
\]
Thus both factors in~\eqref{eq:bloch-factor} are positive, and consequently
$P(x,s)>0$ for every $x>r_*$ and every $s\in[-2,2]$. The polynomial
$P$ is even in $x$, so no eigenvalue has absolute value larger than
$r_*$. Since $z=1$ gives $s=2$ and $r_*$ is a root of the first factor
in \eqref{eq:bloch-factor}, the value $r_*$ is attained. Hence
\[
\rhomax(A_{\sigma^{(8)}})=r_*.
\]
Finally, using $\pi<22/7$ and the alternating Taylor lower bound for cosine gives the
explicit estimates
\[
\cos(\pi/32)>0.9951,
\qquad
\cos(\pi/16)>0.9807.
\]
Therefore
\[
\rho_-(32)
>2\sqrt{0.9951^2+0.9807^2}
>2.7942>2.79361>r_*.
\]
Since both $\cos(\pi/n)$ and $\cos(2\pi/n)$ increase with $n$ for $n\ge32$,
$\rho_-(n)$ is increasing there. Therefore $r_*<\rho_-(n)$ for every $n=8m\ge32$.
\end{proof}

\begin{remark}[Additional periodic search]
As a reproducibility check on Conjecture~\ref{conj:period8}, we also exhaustively scanned
the gauge-fixed subclass in which every step-$1$ edge has sign $+1$ and the step-$2$
signs have period $p\in\{2,4,8,16\}$. For $n=32$, the minimum over all $2^p$
patterns is respectively
$2.828427124746190$ for $p=2$, $2.828427124746190$ for $p=4$,
$2.793604493334841$ for $p=8$, and $2.793604493334841$ for $p=16$.
For $n=64$, the corresponding exhaustive minima for $p=8$ and $p=16$ are again
$2.793604493334841$. These scans are not exhaustive over all signings of $C_n(1,2)$:
they omit, in particular, other periodic gauges and the opposite Hamilton-cycle
holonomy. They are therefore supporting evidence only, and the matching lower bound
in Conjecture~\ref{conj:period8} remains open.
\end{remark}

\subsection{A refined global-optimality conjecture}

\begin{conjecture}[Period-$8$ global optimality]\label{conj:period8}
For every integer $m\ge4$,
\[
\boxed{\min_\sigma\rhomax(A_\sigma)=r_*
\quad\text{on }C_{8m}(1,2),}
\]
where $r_*$ is the largest real root of
\[
x^4-2x^3-6x^2+12x-4=0.
\]
Equivalently, the period-$8$ signing from Theorem~\ref{thm:counterexample}
is a global spectral minimizer on every graph $C_{8m}(1,2)$ with $m\ge4$.
Theorem~\ref{thm:counterexample} proves only the upper bound
$\min_\sigma\rhomax(A_\sigma)\le r_*$; the matching lower bound remains
open.
\end{conjecture}

\begin{remark}[Computational status for the remaining small even sizes]
For $n=26,28,30$, in the heuristic searches described in the
accompanying computational repository, we found no signing with
spectral radius below $\rho_-(n)$. These searches are not exhaustive
proofs of optimality, and we make no claim that $32$ is the first even
value for which the global-optimality statement fails.
\end{remark}

\subsection{\texorpdfstring{The exceptional case $n=8$}{The exceptional case n=8}}

\begin{remark}[The exceptional case $n=8$]\label{rem:n8}
Identifying \eqref{eq:circ-system} with ``every quadrilateral of $C_n(1,2)$ is
unbalanced'' and hence with the parity-family construction of the present paper
requires $n\ge10$. At $n=8$ the step-$2$ edges form two additional
quadrilaterals, $S_0=(0,2,4,6)$ and $S_1=(1,3,5,7)$, which the
canonical signing of Proposition~\ref{prop:circ} balances. All
statements above remain true at $n=8$ for the system \eqref{eq:circ-system}
exactly as written: the proofs of Propositions~\ref{prop:circ}
and~\ref{prop:twist} apply verbatim (rank $7=n-1$, four classes,
$\rho\in\{2\sqrt2,\rho_-(8)\}$, the lattice maximum covered by
$g(\pi/8)>1$). The strictly larger system requiring all ten
quadrilaterals to be unbalanced is resolved by a flux computation.
Over $\F_2$,
\[
\chi_{S_0}=\chi_H+\!\!\sum_{i\ \mathrm{even}}\!\chi_{T_i},\qquad
\chi_{S_1}=\chi_H+\!\!\sum_{i\ \mathrm{odd}}\!\chi_{T_i},
\]
since each sum on the right covers every step-$1$ edge exactly once
and the four step-$2$ edges of the named quadrilateral exactly once.
On the solution family $\tau_i=\tau_0(-1)^i$, so
$\prod_{i\,\mathrm{even}}\tau_i=\prod_{i\,\mathrm{odd}}\tau_i=+1$ and
both extra quadrilaterals carry sign $\alpha$: the two additional
constraints are jointly equivalent to $\alpha=-1$. The all-ten system
therefore has exactly the two twisted classes $(\pm1,-1)$ as its
solutions, both attaining
\[
\rho_-(8)\;=\;2\sqrt{\cos^2(\pi/8)+\cos^2(\pi/4)}\;=\;\sqrt{4+\sqrt2}
\;=\;2.326846\ldots,
\]
and it excludes the $2\sqrt2$ class. The period-$8$ signing of
Theorem~\ref{thm:counterexample} is not a counterexample at $n=8$, since
its spectral radius is approximately $2.7936$, whereas
$\rho_-(8)\approx2.3268$.
\end{remark}

\begin{remark}
The spectral radius is constant on each of the four classes of the
solution family, so the family has only the two values
$\{\rho_-(n),\,2\sqrt2\}$ from Proposition~\ref{prop:twist}. Thus the
alternating triangle flux with anti-periodic step-$1$ holonomy attains
the minimum spectral radius within the quadrilateral-unbalanced family.
Theorem~\ref{thm:counterexample} gives an unrestricted signing with strictly
smaller spectral radius for $n=8m\ge32$. The constrained flux minimization
is nevertheless a useful exact finite model, analogous in spirit to
Lieb's flux-phase viewpoint for signed circulants \cite{lieb94}.


In our sampled instances, uniformly random signings of $C_n(1,2)$ typically had
spectral radius above the Kesten bound for $n\gtrsim480$, while the structured family
remained at $2\sqrt2$ and below. This is an empirical observation from the sampled
instances, not a concentration theorem; it illustrates the distinction between
uniform averaging and structured parity families in the present model.
\end{remark}


\section{Outlook}

The central open problem for the parity-family program is the genuinely
cancellation-sensitive inequality suggested by Proposition~\ref{prop:master}: can the
negative contributions of single-cycle parity classes systematically offset the even-wrap
excess in cycle-dense graphs while the higher-order pair and multiple-cycle classes remain
controlled? The present paper does not prove such a theorem. The bicycle-free result is a
conditioning corollary of known random-signing theory and serves as a baseline for what can
be obtained without exploiting the parity signs quantitatively.

The exact circulant model provides a separate finite laboratory. The twisted class is the
exact minimizer within the quadrilateral-unbalanced family, while the period-$8$ construction
gives a smaller value in the unrestricted problem for every $n=8m\ge32$. Conjecture~\ref{conj:period8}
asks whether this value is the unrestricted minimum on the entire subsequence. The periodic
scan above removes the simplest period-$\le16$ gauge-fixed competitors; a natural next step
is to extend the search to the opposite Hamilton-cycle holonomy and more general periodic
gauges, followed by an analytic lower bound if a stable periodic pattern emerges.
\begingroup
\setlength{\emergencystretch}{2em}
\raggedright
\small\begin{thebibliography}{18}
\bibitem{bass92} H.~Bass, \emph{The Ihara-Selberg zeta function of a
tree lattice}, Internat.~J.~Math.~3 (1992), 717-797.
\bibitem{bilulinial06} Y.~Bilu, N.~Linial,
\emph{Lifts, discrepancy and nearly optimal spectral gap},
Combinatorica 26 (2006), 495-519.
\bibitem{companion} V.~Suvagiya, \emph{Signed circulants at the Ramanujan
bound}, arXiv:2607.18334 (2026).
\bibitem{godsilgutman81} C.~D.~Godsil, I.~Gutman,
\emph{On the matching polynomial of a graph}, in: L.~Lov\'asz and V.~T.~S\'os (eds.),
\emph{Algebraic Methods in Graph Theory} (Proc., Szeged, 1978), Vol.~I,
pp.~241-249, Colloq.~Math.~Soc.~J\'anos~Bolyai, Vol.~25, North-Holland, Amsterdam, 1981.
\bibitem{heilmannlieb72} O.~J.~Heilmann, E.~H.~Lieb,
\emph{Theory of monomer-dimer systems}, Comm.~Math.~Phys.~25 (1972),
190-232.
\bibitem{huang19} H.~Huang, \emph{Induced subgraphs of hypercubes and a
proof of the Sensitivity Conjecture}, Ann.~of Math.~190 (2019), 949-955.
\bibitem{ks00} M.~Kotani, T.~Sunada, \emph{Zeta functions of finite
graphs}, J.~Math.~Sci.~Univ.~Tokyo 7 (2000), 7-25.
\bibitem{lieb94} E.~H.~Lieb, \emph{Flux phase of the half-filled band},
Phys.~Rev.~Lett.~73 (1994), 2158-2161.
\bibitem{mss15} A.~W.~Marcus, D.~A.~Spielman, N.~Srivastava,
\emph{Interlacing families I: Bipartite Ramanujan graphs of all degrees},
Ann.~of Math.~182 (2015), 307-325.
\bibitem{hornjohnson13} R.~A.~Horn, C.~R.~Johnson, \emph{Matrix Analysis}, 2nd ed.,
Cambridge University Press, 2013, Chapter~8.
\bibitem{mop20} S.~Mohanty, R.~O'Donnell, P.~Paredes,
\emph{Explicit near-Ramanujan graphs of every degree}, SIAM J.~Comput.~51
(2022), STOC20-1-STOC20-23; arXiv:1909.06988. DOI: 10.1137/20M1342112.
\bibitem{chenvandambu24} L.~Chen, E.~R.~van Dam, C.~Bu, \emph{Spectra of power
hypergraphs and signed graphs via parity-closed walks}, J. Comb. Theory Ser. A
207 (2024), 105909. DOI: 10.1016/j.jcta.2024.105909.
\bibitem{lihouwang24} D.~Li, Y.~Hou, D.~Wang, \emph{Zeta functions of signed
graphs}, Linear Multilinear Algebra 72 (2024), 2719-2731.
DOI: 10.1080/03081087.2023.2296578.
\bibitem{luoroy24} Y.~Luo, A.~Roy, \emph{Spectral Antisymmetry of Twisted Graph
Adjacency}, arXiv:2403.01550 (2024; revised 2026).
\bibitem{prior26} V.~Suvagiya, \emph{Parity families and a kernel-averaged
$L$-function for near-Ramanujan signings}, arXiv:2607.17343 (2026).
\bibitem{huang26} Z.~Huang, \emph{A $\sqrt2$-Approximation to the Bilu-Linial Conjecture},
arXiv:2609.17551 (2026).
\bibitem{linzhou26} F.~Lin, H.~Zhou, \emph{Improved Bounds for the Bilu-Linial Conjecture via Spectral
Recovery from Mixed Determinantal Polynomials}, arXiv:2609.15715 (2026).
\bibitem{xu26counter} Z.~Xu, \emph{A 3-regular counterexample to the Bilu-Linial signing conjecture},
arXiv:2609.15591 (2026).
\end{thebibliography}
\endgroup

\end{document}
